\documentclass[a4paper,12pt]{book}

% Paquetes necesarios
\usepackage[utf8]{inputenc}   % Codificación de caracteres
\usepackage[spanish]{babel}   % Idioma español
\usepackage[T1]{fontenc}      % Codificación de fuentes
\usepackage{amsmath, amssymb} % Símbolos matemáticos
\usepackage{graphicx}         % Inclusión de gráficos
\usepackage{cite}             % Gestión de citas
\usepackage{hyperref}         % Enlaces y referencias
\usepackage{geometry}         % Configuración de márgenes
\usepackage{fancyhdr}         % Encabezados y pies de página
\usepackage{titlesec}         % Formato de títulos
\usepackage{booktabs}         % Tablas profesionales
\usepackage{caption}          % Personalización de leyendas
\usepackage{enumitem}         % Personalización de listas
\usepackage{float}
\usepackage{tcolorbox}
\usepackage[table]{xcolor} % Paquete para colores en tablas
\usepackage{colortbl}       % Complemento para colorear celdas específicas
\usepackage{multirow}       % Combinar celdas en tablas
\usepackage{makecell}       % Combinar celdas en tablas
\usepackage{enumitem}
\usepackage{amsmath}
\usepackage{eurosym}
\usepackage{tikz}
\usepackage{listings}
\usepackage{color}
\usepackage{float}
\usepackage{pdfpages}
% Configuración de márgenes
\geometry{left=3cm, right=3cm, top=2.5cm, bottom=2.5cm}

% Configuración de encabezados y pies de página
% \setlength{\headheight}{14.49998pt}
\pagestyle{fancy}
\fancyhf{}
\fancyhead[L]{Universidad de Granada}
\fancyhead[L]{\nouppercase{\leftmark}}

% \fancyhead[C]{Escuela Técnica Superior de Ingenierías Informática}
\fancyhead[R]{Fundamentos de Ingeniería del Software}
\fancyfoot[L]{\rule[0pt]{\textwidth}{0.2pt}\\Ismael Sallami Moreno}
\fancyfoot[C]{\rule[0pt]{\textwidth}{0.2pt}\\\thepage}
\fancyfoot[R]{\rule[0pt]{\textwidth}{0.2pt}\\\today}
\renewcommand{\sectionmark}[1]{\markboth{#1}{}} % Configura \leftmark para que solo muestre la sección


% Formato de títulos
\titleformat{\section}{\large\bfseries}{\thesection.}{0.5em}{}
\titleformat{\subsection}{\normalsize\bfseries}{\thesubsection.}{0.5em}{}

% Datos del documento
\title{\textbf{Temario Fundamentos de Ingenierías del Software}}
\author{
    Ismael Sallami Moreno \\
    \texttt{ism350zsallami@correo.ugr.es}
}
\date{
    \vspace{1cm}
    \begin{tabular}{rl}
        \textbf{Asignatura:} & Fundamentos de Base de Datos \\
        \textbf{Tema:} & Teoría \\
        \textbf{Fecha:} & \today
    \end{tabular}
}

\begin{document}

% Portada
\begin{titlepage}
    \begin{center}
        % \vspace*{1cm}
        
        % \Huge
        % \textbf{Práctica Contabilidad Financiera II}
        \Huge \textbf{Temario Fundamentos de Ingeniería del Software} 
        % \vspace{0.5cm}
        % \LARGE
        % \textbf{Ismael Sallami Moreno}\\
        % \LARGE
        % \texttt{ism350zsallami@correo.ugr.es}
        % \LARGE
        % \url{https://github.com/Ismael-Sallami}
        
        % \vfill
        
        % \Large
        % \textbf{Universidad de Granada}
        
        \vspace{0.8cm}
        
        \begin{tikzpicture}[remember picture, overlay]
            \node[opacity=0.2] at (current page.center) {
\includegraphics[width=\paperwidth,height=\paperheight]{portada.jpg}};
            \node[align=center] at (current page.center) {
                
                \vspace{0.5cm}
                \LARGE \textbf{Ismael Sallami Moreno} \\
                \LARGE \texttt{ism350zsallami@correo.ugr.es} \\
                \LARGE \url{https://ismael-sallami.github.io/} \\
                \LARGE \url{https://elblogdeismael.github.io/} \\
                \vspace{2cm}
                \Large \textbf{Universidad de Granada} \\
                \vspace{0.8cm}
                % \Large \textbf{2025}
            };
        \end{tikzpicture}
        \vfill
        
        \Large
        \textbf{2025}
        
    \end{center}
\end{titlepage}
\newpage


\includepdf[pages=-]{licencia.pdf}

% Tabla de contenidos
\tableofcontents
\newpage

% Un capítulo por cada uno de los cuatro temas de la guía docente de la
% asignatura (216113B), en su orden y con sus apartados.
\chapter{Introducción a la Ingeniería del Software}

\section{El producto software, propiedades y ciclo de vida}

El software no se parece a los productos que fabrica el resto de las ingenierías.
No se gasta con el uso, no está sujeto a tolerancias de material y su coste de
replicación es prácticamente nulo: una copia vale lo mismo que un millón. Lo que
cuesta es concebirlo y mantenerlo, y ahí es donde se concentra el esfuerzo de la
disciplina.

De esa naturaleza salen dos consecuencias que condicionan todo lo demás. La
primera es que no hay límites físicos que frenen la complejidad. Un puente no
puede crecer sin que la estructura lo acuse; un sistema software sí, y por eso
puede llegar a un tamaño que ninguna persona abarca. La segunda es que el
software no se deteriora, pero sí se degrada: cada cambio que se le hace lo aleja
del diseño con el que nació, y la facilidad para modificarlo acaba siendo la causa
de que se vuelva difícil de modificar \cite{sommerville2016}.

Conviene separar dos términos que se usan como sinónimos y no lo son. El
\textbf{producto software} es lo que se entrega: los programas, los datos de
configuración y la documentación que permite instalarlo, usarlo y mantenerlo. El
\textbf{sistema software} incluye además el entorno en el que ese producto opera.
Entregar solo los ejecutables no es entregar el producto.

\subsection*{Propiedades de un producto software}

La calidad de un producto software no se mide únicamente por que haga lo que se le
pide. Las propiedades que suelen exigirse son:

\begin{itemize}
    \item \textbf{Corrección y fiabilidad.} Que haga lo especificado y que siga
          haciéndolo en condiciones adversas, incluidas las que no se previeron.
    \item \textbf{Mantenibilidad.} Que pueda evolucionar con un coste razonable.
          Es la propiedad que más peso tiene a largo plazo, porque el grueso del
          gasto de un sistema se produce después de la primera entrega.
    \item \textbf{Eficiencia.} Que use de forma razonable el tiempo de cómputo, la
          memoria y el resto de recursos.
    \item \textbf{Usabilidad.} Que las personas a las que va dirigido puedan
          usarlo sin formación desproporcionada.
    \item \textbf{Seguridad y protección.} Que no cause daños y que resista un uso
          malintencionado.
\end{itemize}

Estas propiedades compiten entre sí. Un diseño muy eficiente suele ser menos
legible y, por tanto, menos mantenible; una interfaz muy flexible complica la
validación. Parte del trabajo de ingeniería consiste precisamente en decidir qué
se sacrifica y dejar constancia de por qué.

\subsection*{El ciclo de vida}

El ciclo de vida es el conjunto de etapas por las que pasa un producto desde que
se identifica la necesidad hasta que se retira. Con distintos nombres, casi todos
los modelos reconocen las mismas actividades: especificación, diseño e
implementación, validación y evolución.

Lo que distingue a un modelo de otro no son las actividades, sino cómo se
ordenan y con qué frecuencia se repiten. Merece la pena insistir en la última
etapa: \textbf{la evolución no es un apéndice del ciclo de vida, es la mayor parte
de él}. Un sistema en producción recibe cambios durante años, y el diseño con el
que se construyó determina lo caros que resultan.

\section{Concepto de Ingeniería del Software}

El término nació en la conferencia de la OTAN de 1968, convocada porque los
proyectos software se entregaban tarde, por encima del presupuesto y sin la
calidad prometida. La expresión «crisis del software» describía esa situación, y
la propuesta fue tratar el desarrollo como una ingeniería: con métodos, medidas y
procesos, en lugar de como un oficio artesanal.

La Ingeniería del Software es, en esa línea, la disciplina que se ocupa de todos
los aspectos de la producción de software, desde las primeras etapas de la
especificación hasta el mantenimiento posterior a la entrega
\cite{sommerville2016}. Dos matices de esa definición importan:

\begin{itemize}
    \item \textbf{Todos los aspectos}, no solo la programación. La gestión del
          proyecto, la elección de herramientas y la organización del equipo
          forman parte de la disciplina.
    \item \textbf{Hasta el mantenimiento}, no hasta la entrega. Un método que
          funciona bien hasta el día del despliegue y deja el sistema
          intratable después no ha resuelto el problema.
\end{itemize}

Conviene distinguirla de dos disciplinas vecinas. La \textbf{informática} estudia
los fundamentos teóricos del cómputo; la ingeniería del software se ocupa de
construir sistemas útiles con esos fundamentos, aceptando restricciones de plazo,
coste y personas. La \textbf{ingeniería de sistemas} es más amplia: abarca el
hardware, los procesos y la organización, y el software es solo una de sus
partes.

\subsection*{Responsabilidad profesional}

La disciplina lleva asociada una dimensión ética que no es decorativa. Quien
construye software decide, en la práctica, sobre la privacidad de terceros, sobre
la accesibilidad de un servicio y, en sistemas críticos, sobre la seguridad de
las personas. Los códigos deontológicos de ACM e IEEE recogen los compromisos
habituales: confidencialidad, competencia declarada con honestidad, respeto a la
propiedad intelectual y no aceptar encargos que se sepan perjudiciales
\cite{sommerville2016}.

\section{El proceso de desarrollo de software}

Un proceso de desarrollo es un conjunto ordenado de actividades que produce un
producto software. Los modelos que se estudian a continuación no son recetas
excluyentes: en un proyecto real se combinan, y la elección depende de cuánto se
conoce el problema al empezar y de cuánto cuesta equivocarse.

\subsection*{Modelo en cascada}

Ordena las actividades en fases sucesivas —especificación, diseño,
implementación, verificación y mantenimiento— y exige cerrar cada una antes de
abrir la siguiente. Su ventaja es la trazabilidad: cada fase deja documentos
revisables, lo que encaja bien en desarrollos con requisitos contractuales o
sujetos a certificación.

Su inconveniente es el que lo hizo célebre: los requisitos rara vez se conocen por
completo al principio, y el modelo hace caro cambiarlos. Un error de
especificación no se detecta hasta la validación, cuando corregirlo cuesta
órdenes de magnitud más que en su fase de origen \cite{pressman2021}.

\subsection*{Desarrollo incremental}

Divide el sistema en incrementos que se especifican, construyen y entregan por
separado, empezando por los de mayor valor. Cada entrega es funcional y sirve para
obtener realimentación antes de comprometerse con el resto.

Reduce el coste de acomodar cambios y adelanta el momento en que el cliente ve
algo que funciona. A cambio, el proceso es menos visible desde fuera —no hay
documentos de fase que marquen el avance— y la estructura del sistema tiende a
degradarse si no se reserva esfuerzo para reorganizarla.

\subsection*{Modelos evolutivos y por prototipos}

Cuando el problema no está claro, resulta razonable construir una versión reducida
que sirva para entenderlo. El prototipo puede ser \emph{desechable}, si su único
fin es aclarar requisitos y luego se descarta, o \emph{evolutivo}, si se refina
hasta convertirse en el producto. El riesgo del segundo es conocido: un prototipo
se construye sin las garantías que se exigen a un producto, y promoverlo tal cual
deja una base frágil.

El \textbf{modelo en espiral} de Boehm organiza el desarrollo en ciclos y sitúa el
análisis de riesgos como actividad explícita al principio de cada uno. Su
aportación duradera no es la forma de la espiral, sino la idea de que el proceso
debe adaptarse al riesgo asumido.

\subsection*{Métodos ágiles}

Los métodos ágiles parten de una constatación: en muchos proyectos el coste de
predecir es mayor que el de adaptarse. Priorizan las entregas frecuentes, la
colaboración con el cliente y la respuesta al cambio por encima de la planificación
detallada y la documentación exhaustiva. Scrum organiza el trabajo en iteraciones
cortas con roles y reuniones definidos; la programación extrema añade prácticas
técnicas como la integración continua, la programación por parejas y el desarrollo
guiado por pruebas \cite{beck2002,lasa2017}.

No son la respuesta a todo. Su eficacia se apoya en equipos pequeños, en un cliente
disponible y en un dominio donde equivocarse sea barato. En sistemas con
certificación obligatoria, con muchos equipos coordinados o con un coste de fallo
muy alto, conviene un proceso más ceremonioso, aunque adopte prácticas ágiles en
lo técnico.

\chapter{Ingeniería de requisitos}

\section{Introducción a la ingeniería de requisitos}

Determinar qué debe hacer un sistema es más difícil que construirlo. El problema
no es técnico: consiste en enunciar con precisión un problema complejo del que
las personas implicadas tienen ideas distintas, a veces incompatibles y casi
siempre incompletas.

Un \textbf{requisito} es una condición o capacidad que el sistema debe satisfacer.
La palabra se usa con dos niveles de detalle que conviene no mezclar
\cite{sommerville2016}:

\begin{itemize}
    \item Los \textbf{requisitos de usuario} se enuncian en lenguaje natural,
          desde el punto de vista de quien va a usar el sistema, y sirven para
          acordar el alcance con el cliente.
    \item Los \textbf{requisitos del sistema} detallan las funciones, servicios y
          restricciones con la precisión necesaria para diseñar y contratar.
\end{itemize}

La distinción clásica es otra, y es transversal a la anterior:

\begin{itemize}
    \item \textbf{Requisitos funcionales}: lo que el sistema debe hacer.
    \item \textbf{Requisitos no funcionales}: las restricciones sobre cómo debe
          hacerlo —rendimiento, disponibilidad, seguridad, normativa aplicable—.
\end{itemize}

Los no funcionales suelen recibir menos atención y son los que más condicionan la
arquitectura. Un requisito de tiempo de respuesta o de tolerancia a fallos afecta
al sistema entero, mientras que un requisito funcional suele quedar confinado a
un módulo. Además, incumplir un requisito no funcional puede inutilizar el sistema
aunque todas sus funciones estén implementadas.

\subsection*{Propiedades de un buen conjunto de requisitos}

Un requisito bien escrito es \emph{necesario}, \emph{verificable}, \emph{no
ambiguo} y \emph{consistente} con los demás. La verificabilidad es la propiedad
que más se descuida: «el sistema debe ser rápido» no es un requisito, porque no
existe prueba que decida si se cumple. «El sistema debe responder a una consulta
de catálogo en menos de dos segundos con cien usuarios concurrentes» sí lo es.

Al conjunto se le exige además que sea \emph{completo} y \emph{trazable}: que
cubra todo lo necesario y que cada elemento del diseño y de las pruebas pueda
remontarse al requisito que lo justifica. La completitud absoluta es inalcanzable
en sistemas grandes, y esa es una de las razones por las que los procesos
iterativos ganaron terreno.

\subsection*{Delimitar lo que queda fuera}

Tan importante como decir qué hace el sistema es dejar por escrito qué no hace. Un
alcance que no se cierra crece durante el desarrollo sin que nadie lo decida, y
ese crecimiento es una de las causas habituales de que un proyecto se desvíe en
plazo y coste. La lista de exclusiones debe ser explícita y estar acordada.

\section{Obtención de requisitos}

La obtención —o elicitación— es la actividad de descubrir los requisitos hablando
con las personas implicadas y observando cómo trabajan. Las dificultades son
recurrentes: quienes van a usar el sistema describen su trabajo en términos de la
herramienta que ya conocen, dan por supuesto lo que para ellos es evidente, y
distintos grupos tienen objetivos que chocan.

Las técnicas más usadas son:

\begin{itemize}
    \item \textbf{Entrevistas}, abiertas o estructuradas. Son el instrumento
          principal. Las abiertas sirven para explorar; las estructuradas, para
          confirmar y precisar.
    \item \textbf{Observación} del trabajo real. Descubre lo que nadie menciona en
          una entrevista porque lo hace de forma automática, y también los atajos
          con los que se sortea el sistema actual.
    \item \textbf{Talleres} con varios grupos a la vez, útiles para sacar a la luz
          los conflictos entre ellos en vez de descubrirlos al final.
    \item \textbf{Prototipos} y maquetas, que convierten una discusión abstracta
          en una reacción concreta ante algo visible.
    \item \textbf{Estudio de la documentación} existente y del sistema al que se
          sustituye.
\end{itemize}

Identificar a los \emph{stakeholders} —toda persona u organización afectada por el
sistema— forma parte de la actividad. Olvidar a un grupo es la vía más directa a
un requisito descubierto tarde: los administradores del sistema, el personal de
soporte y los responsables de cumplimiento normativo quedan fuera de la lista con
frecuencia y aportan restricciones que luego resultan innegociables.

\section{Modelado de casos de uso e historias de usuario}

\subsection*{Casos de uso}

Un caso de uso describe una interacción entre el sistema y un actor externo que
produce un resultado con valor para ese actor. El \textbf{actor} es un rol, no una
persona: el mismo empleado puede actuar como «cliente» y como «administrador» en
casos distintos.

El diagrama de casos de uso de UML solo aporta la vista general: qué actores hay y
con qué casos se relacionan. El contenido está en la descripción textual de cada
caso, que suele incluir el actor principal, las precondiciones, el flujo
principal paso a paso, los flujos alternativos y la postcondición
\cite{arlow2006,larman2006}.

Las relaciones entre casos —\texttt{include} para el comportamiento común y
\texttt{extend} para el opcional— conviene usarlas con moderación. Un modelo lleno
de descomposiciones se vuelve ilegible y empieza a describir el diseño en vez del
comportamiento observable, que es justo lo que un caso de uso debe evitar.

El error más frecuente al modelar es escribir los casos desde dentro del sistema,
en términos de pantallas y de operaciones sobre la base de datos. Un caso de uso
debe poder leerlo el cliente y reconocer su trabajo en él.

\subsection*{Historias de usuario}

Las historias de usuario son la alternativa ágil: enunciados breves, en la forma
«como \emph{rol} quiero \emph{acción} para \emph{beneficio}», acompañados de
criterios de aceptación. No pretenden documentar el sistema, sino recordar una
conversación pendiente con el cliente y servir de unidad de planificación.

Frente a los casos de uso son más ligeras y más fáciles de reordenar por
prioridad, pero dejan menos rastro. La elección entre unos y otras no es de moda:
depende de si el cliente está disponible de forma continua —caso en el que la
conversación sustituye al documento— o de si el acuerdo tiene que sostenerse por
sí solo durante meses.

\section{Especificación y análisis}

\subsection*{Especificación}

Especificar es dejar los requisitos por escrito de forma que sirvan de contrato
técnico. El documento de especificación —el \emph{SRS}, por sus siglas en
inglés— recoge la introducción y el alcance, la descripción general del sistema y
su entorno, los requisitos funcionales, los no funcionales y las restricciones,
más los apéndices con glosario y modelos.

El estilo importa. Los requisitos se enuncian en frases cortas, una condición por
requisito, con identificador propio y sin conjunciones que escondan dos
exigencias en una. Las palabras que expresan obligación deben usarse con un
significado acordado y constante, y conviene evitar los términos que no admiten
prueba: «adecuado», «amigable», «eficiente», «si es posible».

\subsection*{Análisis}

El análisis transforma los requisitos en modelos que permiten razonar sobre ellos
antes de diseñar. En un enfoque orientado a objetos, el resultado típico es el
\textbf{modelo conceptual} o modelo del dominio: las clases del problema, sus
atributos y las asociaciones entre ellas \cite{larman2006}.

Ese modelo no es un diseño. Recoge los conceptos del dominio tal como existen para
el cliente, sin métodos, sin decisiones de implementación y sin clases técnicas.
Confundir ambas cosas produce un modelo del dominio contaminado de controladores y
gestores, que ya no sirve para hablar con quien conoce el negocio.

Junto al modelo conceptual se emplean los \textbf{diagramas de secuencia del
sistema}, que muestran las operaciones que un actor invoca sobre el sistema visto
como una caja negra, y los \textbf{contratos de operación}, que precisan el efecto
de cada una en términos de precondiciones y postcondiciones.

\subsection*{Validación y gestión de los requisitos}

Validar los requisitos es comprobar que describen el sistema que el cliente quiere,
frente a verificar, que es comprobar que el sistema construido cumple los
requisitos. Las técnicas habituales son las revisiones formales con los
implicados, la construcción de prototipos y la redacción anticipada de los casos
de prueba: un requisito para el que no se sabe escribir una prueba está mal
enunciado.

Los requisitos cambian durante todo el proyecto, así que necesitan gestión: una
línea base acordada, un procedimiento para proponer y aprobar cambios, y
trazabilidad para saber qué se ve afectado por cada uno. Sin trazabilidad, el
coste de un cambio se estima a ojo y siempre por debajo.

\chapter{Diseño del software}

\section{Conceptos y principios de diseño}

El diseño es la actividad que decide cómo se organiza el sistema para satisfacer
los requisitos. Entre la especificación y el código hay un espacio de decisiones
que no está determinado por ninguno de los dos, y es ahí donde se juega la
mantenibilidad futura.

Los conceptos que sostienen esas decisiones son pocos y llevan décadas siendo los
mismos \cite{pressman2021}:

\begin{itemize}
    \item \textbf{Abstracción.} Trabajar con una descripción que retiene lo
          relevante y omite el resto, y hacerlo por niveles.
    \item \textbf{Modularidad.} Dividir el sistema en partes con una
          responsabilidad identificable, de forma que cada una pueda entenderse y
          modificarse por separado.
    \item \textbf{Ocultación de información.} Que cada módulo esconda sus
          decisiones internas tras una interfaz. Es lo que permite cambiar la
          implementación sin arrastrar a los clientes del módulo.
    \item \textbf{Separación de intereses.} Que cada parte se ocupe de una
          cuestión y no de varias entremezcladas.
\end{itemize}

De ellos se derivan las dos medidas con las que se juzga una descomposición:

\begin{itemize}
    \item La \textbf{cohesión} mide hasta qué punto los elementos de un módulo
          contribuyen a un mismo propósito. Interesa que sea alta.
    \item El \textbf{acoplamiento} mide la dependencia entre módulos. Interesa que
          sea bajo, y sobre todo que las dependencias vayan hacia abstracciones
          estables y no hacia detalles.
\end{itemize}

Las dos van juntas: una descomposición que reparte una misma responsabilidad entre
varios módulos baja la cohesión y sube el acoplamiento a la vez, porque esos
módulos tendrán que cambiar siempre juntos.

Por encima de estas medidas están los principios que guían la asignación de
responsabilidades. Los patrones GRASP —experto en información, creador,
controlador, bajo acoplamiento, alta cohesión— dan un vocabulario para justificar
por qué una operación va en una clase y no en otra, que es la decisión más
frecuente del diseño orientado a objetos \cite{larman2006}.

\section{Diseño de los casos de uso}

El diseño de casos de uso conecta el análisis con la estructura de objetos:
partiendo de las operaciones del sistema identificadas en el análisis, se decide
qué objetos colaboran para realizarlas y con qué mensajes.

El instrumento habitual es el \textbf{diagrama de interacción}, en cualquiera de
sus dos formas. El diagrama de secuencia ordena los mensajes en el tiempo y se lee
bien cuando el flujo importa; el diagrama de comunicación destaca los enlaces
entre objetos y se lee mejor cuando lo que interesa es la estructura de
colaboración \cite{arlow2006}.

Diseñar una interacción obliga a repartir responsabilidades, y ese reparto es el
que determina el acoplamiento del resultado. Dos criterios ayudan a decidirlo:

\begin{itemize}
    \item Asignar cada responsabilidad a la clase que ya tiene la información
          necesaria para cumplirla, en vez de crear una clase que pida esa
          información a todas las demás.
    \item Introducir un \emph{controlador} que reciba los eventos del sistema y
          delegue, en lugar de dejar que la capa de presentación hable
          directamente con el dominio.
\end{itemize}

Cuando un objeto necesita datos de muchos otros para hacer su trabajo, la
interacción lo delata: el diagrama se llena de mensajes de consulta salientes.
Esa forma es la señal de que la responsabilidad está mal colocada.

\section{Diseño de la estructura de objetos}

El resultado del diseño de las interacciones se consolida en el \textbf{diagrama
de clases de diseño}, que añade al modelo conceptual lo que el análisis
deliberadamente omitía: las operaciones, la visibilidad de los atributos, los
tipos, la navegabilidad de las asociaciones y las clases que no pertenecen al
dominio sino a la solución.

Hay que distinguirlo del modelo conceptual del que procede. Aquel describía el
problema; este describe la solución, y por eso puede contener clases que no
existen en el dominio —repositorios, adaptadores, fábricas— sin que eso sea un
defecto.

Sobre esa estructura se aplican los \textbf{patrones de diseño}, soluciones
probadas a problemas recurrentes. Los que más se emplean en este nivel son los de
creación, para desacoplar al cliente de la clase concreta que instancia; los
estructurales, como adaptador o fachada, para acomodar interfaces que no encajan;
y los de comportamiento, como estrategia u observador, para dejar variable lo que
se sabe que va a cambiar \cite{larman2006}.

Los patrones se aplican cuando el problema que resuelven está presente. Aplicarlos
por anticipado, «por si acaso», añade indirección sin beneficio y es una forma
frecuente de complicar un diseño sin necesidad.

\section{Arquitectura del software}

La arquitectura es el conjunto de decisiones de diseño de mayor alcance: la
descomposición del sistema en partes principales, la relación entre ellas y las
propiedades globales que esa organización garantiza. Son las decisiones más caras
de revertir, porque el resto del diseño se apoya en ellas.

La arquitectura la determinan, sobre todo, los requisitos no funcionales. El
rendimiento empuja hacia componentes gruesos y poca comunicación entre ellos; la
seguridad, hacia capas con control de acceso en las fronteras; la disponibilidad,
hacia la redundancia; la mantenibilidad, hacia componentes pequeños y
reemplazables. Como estas fuerzas tiran en direcciones opuestas, no existe una
arquitectura buena en abstracto, solo una adecuada a unas prioridades declaradas
\cite{sommerville2016}.

Los estilos más comunes son:

\begin{itemize}
    \item \textbf{En capas.} Cada capa usa solo servicios de la inmediatamente
          inferior. Facilita sustituir una capa entera y aísla el dominio de la
          tecnología, a costa de rendimiento en las llamadas que la atraviesan.
    \item \textbf{Cliente-servidor.} Separa a quien pide de quien sirve, lo que
          permite escalar cada lado por separado.
    \item \textbf{Modelo-vista-controlador.} Aísla la presentación del estado y de
          la lógica, de modo que varias vistas puedan compartir un mismo modelo.
    \item \textbf{Por componentes o servicios.} Divide el sistema en unidades
          desplegables de forma independiente. Gana en escalabilidad y en
          autonomía de los equipos; paga con complejidad de operación y con los
          problemas propios de la comunicación en red.
    \item \textbf{Tuberías y filtros.} Encadena transformaciones sobre un flujo de
          datos. Encaja bien en procesamiento por lotes.
\end{itemize}

Decidir la arquitectura incluye documentarla. Una arquitectura que solo existe en
la cabeza de quien la diseñó deja de cumplirse en cuanto el equipo crece, y la
erosión resultante es difícil de revertir porque el código ya depende de la forma
degradada.

\chapter{Otros aspectos de la Ingeniería del Software}

\section{Planificación y gestión de proyectos software}

Gestionar un proyecto software consiste en conseguir que se entregue en plazo,
dentro del presupuesto y con la calidad acordada. Es una actividad de ingeniería
tanto como el diseño, y comparte con él la dificultad de decidir con información
incompleta.

Los proyectos software tienen rasgos que los distinguen de otros: el producto es
intangible, así que el avance no se ve; no hay dos proyectos iguales, de modo que
la experiencia previa se traslada mal; y el proceso cambia con la tecnología a un
ritmo que envejece rápido los datos históricos \cite{sommerville2016}.

\subsection*{Estimación}

Estimar es predecir esfuerzo, plazo y coste. Los enfoques habituales son tres, y
suelen combinarse:

\begin{itemize}
    \item \textbf{Juicio de expertos}, con varias personas estimando por separado
          y contrastando después. Es barato y sorprendentemente robusto cuando el
          equipo conoce el dominio.
    \item \textbf{Estimación por analogía}, comparando con proyectos anteriores
          medidos.
    \item \textbf{Modelos algorítmicos}, como COCOMO, que calculan el esfuerzo a
          partir de una medida del tamaño y de factores de ajuste.
\end{itemize}

Todos comparten el mismo punto débil: dependen de una estimación del tamaño que se
hace cuando menos se sabe del sistema. De ahí la práctica de estimar por rangos,
revisar la estimación en cada iteración y desconfiar de las cifras cerradas dadas
antes de la especificación.

Conviene recordar además que el esfuerzo y el plazo no son intercambiables.
Añadir personas a un proyecto retrasado aumenta el coste de coordinación y suele
retrasarlo más, observación que Brooks formuló hace medio siglo y que sigue
describiendo bien lo que ocurre.

\subsection*{Planificación y riesgos}

La planificación descompone el trabajo en tareas con dependencias, asigna
recursos y fija hitos. Los hitos deben corresponder a entregables verificables;
un hito que se declara cumplido porque «está casi terminado» no informa de nada.

La gestión de riesgos es la parte de la planificación que se salta con más
frecuencia y la que más pronto se echa en falta. Consiste en identificar lo que
puede salir mal, estimar su probabilidad y su impacto, decidir qué hacer con cada
riesgo —evitarlo, mitigarlo, transferirlo o aceptarlo— y revisar la lista de forma
periódica. Los riesgos habituales en software son la rotación del personal, el
cambio de requisitos, la subestimación del tamaño y la dependencia de tecnología
o proveedores no probados.

\subsection*{Calidad y métricas}

La gestión de la calidad establece qué se entiende por calidad en el proyecto y
cómo se comprueba: estándares de producto, estándares de proceso, revisiones e
inspecciones. Las métricas —tamaño, complejidad, densidad de defectos, cobertura
de pruebas— sirven para tomar decisiones, no para evaluar personas. Usadas como
objetivo, dejan de medir lo que medían, porque el equipo optimiza el indicador y
no la propiedad que representaba.

\section{Validación y verificación de software}

Verificar es comprobar que el sistema se construye conforme a su especificación;
validar es comprobar que la especificación describe lo que el cliente necesita. La
distinción clásica lo resume bien: la verificación pregunta si se está
construyendo el producto correctamente, y la validación, si se está construyendo
el producto correcto.

\subsection*{Pruebas}

Las pruebas ejecutan el sistema con la intención de encontrar defectos. Un
resultado sin fallos no demuestra que no los haya: las pruebas muestran la
presencia de errores, nunca su ausencia. Se organizan en niveles:

\begin{itemize}
    \item \textbf{Unitarias}, sobre componentes aislados, escritas normalmente por
          quien los implementa y automatizadas.
    \item \textbf{De integración}, sobre las interacciones entre componentes, donde
          aparecen los fallos de interfaz que las unitarias no ven.
    \item \textbf{De sistema}, sobre el producto completo, incluidos los requisitos
          no funcionales: carga, seguridad, recuperación.
    \item \textbf{De aceptación}, ejecutadas con el cliente o en sus términos, y
          que pertenecen a la validación más que a la verificación.
\end{itemize}

El diseño de los casos de prueba puede partir de la especificación —prueba de caja
negra, con técnicas como las clases de equivalencia y el análisis de valores
límite— o de la estructura interna —prueba de caja blanca, guiada por criterios de
cobertura—. Ninguno de los dos basta por sí solo: la caja negra no llega a los
caminos no previstos y la caja blanca no detecta lo que falta implementar
\cite{pfleeger2002}.

La \textbf{prueba de regresión} es la que se vuelve a pasar tras cada cambio para
comprobar que no se ha roto lo que funcionaba. Solo es viable si está
automatizada, y por eso la automatización de pruebas es una decisión de proceso
más que una preferencia técnica. En la programación extrema se lleva al extremo
escribiendo la prueba antes que el código \cite{beck2002}.

\subsection*{Revisiones e inspecciones}

No toda verificación requiere ejecutar. Las revisiones e inspecciones examinan
documentos y código con una lista de comprobación, y encuentran defectos que las
pruebas no alcanzan, entre ellos los de la propia especificación. Su ventaja es
que actúan antes: un defecto detectado en la revisión de requisitos cuesta una
fracción de lo que cuesta corregido en producción.

\section{Mantenimiento de software}

El mantenimiento es todo el trabajo que se hace sobre el sistema después de la
entrega, y es la fase más larga y más cara del ciclo de vida. Se clasifica en
cuatro tipos:

\begin{itemize}
    \item \textbf{Correctivo}, para reparar defectos.
    \item \textbf{Adaptativo}, para acomodar cambios del entorno: sistema
          operativo, normativa, sistemas con los que se integra.
    \item \textbf{Perfectivo}, para añadir funcionalidad o mejorar el rendimiento a
          petición del usuario.
    \item \textbf{Preventivo}, para facilitar el mantenimiento futuro sin cambiar
          el comportamiento observable.
\end{itemize}

Contra la intuición, el correctivo es la parte menor. La mayoría del esfuerzo se
va en perfectivo y adaptativo, es decir, en hacer evolucionar un sistema que
funciona \cite{sommerville2016}.

\subsection*{Evolución y deterioro}

Las leyes de Lehman describen lo que le ocurre a un sistema en uso: si no cambia,
deja de resultar útil, porque el entorno sí cambia; y a medida que cambia, su
complejidad crece salvo que se dedique trabajo explícito a contenerla. De ahí que
el mantenimiento preventivo —refactorizar, actualizar dependencias, recuperar la
documentación— no sea un lujo, sino la condición para que los otros tres tipos
sigan siendo asequibles.

La \textbf{deuda técnica} nombra ese fenómeno desde el lado económico: cada
decisión tomada por rapidez en lugar de por calidad genera un interés que se paga
en cada cambio posterior. Es legítimo contraerla de forma consciente y con un plan
para devolverla; el problema aparece cuando se contrae sin registrarla, porque
entonces nadie decide cuándo se paga.

\subsection*{Reingeniería}

Cuando el coste de mantener supera al de rehacer, cabe la reingeniería: reconstruir
el sistema conservando su funcionalidad. Es menos arriesgada que un desarrollo
desde cero, porque el sistema existente documenta el comportamiento esperado,
incluido el que nadie llegó a escribir en ningún requisito. La alternativa de
reescribir de cero pierde ese conocimiento y suele reintroducir defectos que ya
estaban resueltos.


\newpage
% Las obras de la bibliografía fundamental y complementaria de la guía docente
% que este temario usa de verdad. No se listan las que no se citan.
\begin{thebibliography}{99}

\bibitem{sommerville2016}
I. Sommerville, \emph{Software Engineering}, 10\textsuperscript{a} ed., Pearson,
2016.

\bibitem{pressman2021}
R. S. Pressman, \emph{Ingeniería del Software: un enfoque práctico}, McGraw-Hill,
2021.

\bibitem{larman2006}
C. Larman, \emph{UML y patrones: una introducción al análisis y diseño orientado
a objetos y al proceso unificado}, 3\textsuperscript{a} ed., Pearson, 2006.

\bibitem{arlow2006}
J. Arlow e I. Neustadt, \emph{UML 2}, Anaya Multimedia, 2006.

\bibitem{pfleeger2002}
S. L. Pfleeger, \emph{Ingeniería de Software: teoría y práctica}, Prentice Hall,
2002.

\bibitem{lasa2017}
C. Lasa Gómez, A. Álvarez García y R. de las Heras del Dedo, \emph{Métodos
Ágiles: Scrum, Kanban, Lean}, Anaya Multimedia, 2017.

\bibitem{beck2002}
K. Beck, \emph{Una explicación de la Programación Extrema: aceptar el cambio},
Addison-Wesley, 2002.

\end{thebibliography}

\end{document}
